% !TEX root =./ECE448_Practice_main.tex
%
% L19 practice set (AIS). Every question is a step in the L19 flowgraphs:
% slots and bursts, bit stuffing and NRZI by hand, the receiver's settings for a
% recording, reading an NMEA sentence, and explaining what noise and the squelch do.
% Blank copy: \begin{solution}[h] reserves h of work space.
% SOLUTIONS copy (\iskey=1) prints the red worked solutions.

\fullwidth{\textbf{Learning Objective 6.8: } I can receive and decode AIS messages from
marine vessels.
\\[4pt]\normalfont\small Show your work. Write a \textbf{Documentation} line at the
end (\textbf{None} if you did not collaborate).}

\begin{questions}

%%%%%%%%%%%%%%%%%%%%%%%%%%%%%%%%%%%%%%%%%
\question \textbf{Slots.} Each AIS channel divides every minute into 2,250 slots, and
sends 9600 b/s.
\begin{parts}

	\begin{minipage}[t]{0.58\linewidth}
		\part Find the length of one slot in milliseconds and in bits.
		\begin{solution}[1.0in]
		\eq{ \frac{60\ \text{s}}{2250} = 26.67\ \text{ms}, \qquad
		9600 \times \frac{60}{2250} = 256\ \text{bits} }
		\end{solution}
	\end{minipage}\hfill%
	\begin{minipage}[t]{0.37\linewidth}
		\raggedleft \vspace{12pt}
		ms: \ansbox[0.9in]{26.67}

		\vspace{4pt}
		bits: \ansbox[0.9in]{256}
	\end{minipage}

	\begin{minipage}[t]{0.58\linewidth}
		\part A burst carries an 8-bit ramp, 24 training bits, two 8-bit flags, the message,
		and a 16-bit CRC. How many slots does a 168-bit position report need, and how
		many does a 424-bit name and voyage report need, before stuffed bits?
		\begin{solution}[1.1in]
		The overhead is $8 + 24 + 8 + 16 + 8 = 64$ bits. A position report is
		$64 + 168 = 232$ bits, which fits in one slot of 256 with room for stuffed bits.
		A name and voyage report is $64 + 424 = 488$ bits, so it needs two slots.
		\end{solution}
	\end{minipage}\hfill%
	\begin{minipage}[t]{0.37\linewidth}
		\raggedleft \vspace{12pt}
		position: \ansbox[0.9in]{1}

		\vspace{4pt}
		name: \ansbox[0.9in]{2}
	\end{minipage}

\end{parts}

\newpage
%%%%%%%%%%%%%%%%%%%%%%%%%%%%%%%%%%%%%%%%%
\question \textbf{Stuffing and NRZI by hand.} The data bits \texttt{011111111100} are
about to be sent. The transmitter stuffs a 0 after every five 1s in a row, then applies
NRZI: a 0 changes the level and a 1 keeps it. The level before the first bit is high (H).
\begin{parts}

	\begin{minipage}[t]{0.58\linewidth}
		\part Write the bits after stuffing.
		\begin{solution}[0.9in]
		After the first five 1s a 0 goes in; the count then restarts, and the four 1s that
		remain are fewer than five: \texttt{0111110111100}, 13 bits.
		\end{solution}
	\end{minipage}\hfill%
	\begin{minipage}[t]{0.37\linewidth}
		\raggedleft \vspace{12pt}
		\ansbox[1.5in]{0111110111100}
	\end{minipage}

	\begin{minipage}[t]{0.58\linewidth}
		\part Write the level, H or L, of each of the 13 transmitted bits.
		\begin{solution}[1.0in]
		The first 0 changes H to L, the five 1s keep L, the stuffed 0 changes to H, the
		four 1s keep H, and the last two 0s change to L and back to H:
		\texttt{LLLLLLHHHHHLH}.
		\end{solution}
	\end{minipage}\hfill%
	\begin{minipage}[t]{0.37\linewidth}
		\raggedleft \vspace{12pt}
		\ansbox[1.5in]{LLLLLLHHHHHLH}
	\end{minipage}

	\part In words, explain why the receiver does not need to know which FM frequency,
	up or down, stands for H.
		\begin{solution}[1.0in]
		Under NRZI the bit is in whether the level changed, not in which level it is. If
		the receiver swapped H and L, every change would still be a change and every
		repeat a repeat, so the bits come out the same.
		\end{solution}

\end{parts}

\newpage
%%%%%%%%%%%%%%%%%%%%%%%%%%%%%%%%%%%%%%%%%
\question \textbf{Receiver settings.} A classmate brings back a SigMF recording whose
\texttt{.sigmf-meta} file says \texttt{cf32\_le}, a sample rate of 250000, and a
center frequency of 162000000.
\begin{parts}

	\begin{minipage}[t]{0.58\linewidth}
		\part Set the Frequency Xlating FIR Filter's center frequency for channel A and
		for channel B.
		\begin{solution}[0.9in]
		Channel A is $161.975 - 162.000 = -0.025$ MHz, or $-25$ kHz; channel B is
		$162.025 - 162.000 = +25$ kHz.
		\end{solution}
	\end{minipage}\hfill%
	\begin{minipage}[t]{0.37\linewidth}
		\raggedleft \vspace{12pt}
		A: \ansbox[1.0in]{$-25$ kHz}

		\vspace{4pt}
		B: \ansbox[1.0in]{$+25$ kHz}
	\end{minipage}

	\begin{minipage}[t]{0.58\linewidth}
		\part Choose \texttt{decim} so that the rate after the filter is about 48 kS/s, and
		find the samples per bit that Symbol Sync needs.
		\begin{solution}[1.0in]
		$250/48 \approx 5.2$, so \texttt{decim} $= 5$, which gives 50 kS/s and
		$50000/9600 = 5.21$ samples per bit. Symbol Sync accepts a fractional value.
		\end{solution}
	\end{minipage}\hfill%
	\begin{minipage}[t]{0.37\linewidth}
		\raggedleft \vspace{12pt}
		decim: \ansbox[0.9in]{5}

		\vspace{4pt}
		sps: \ansbox[0.9in]{5.21}
	\end{minipage}

	\begin{minipage}[t]{0.58\linewidth}
		\part The Quadrature Demod's gain is the rate after decimation divided by
		$2\pi \times 2400$, so that a 2.4 kHz shift comes out as 1. Find it.
		\begin{solution}[0.9in]
		\eq{ \frac{50000}{2\pi \times 2400} = 3.32 }
		\end{solution}
	\end{minipage}\hfill%
	\begin{minipage}[t]{0.37\linewidth}
		\raggedleft \vspace{12pt}
		\ansbox[0.9in]{3.32}
	\end{minipage}

\end{parts}

\newpage
%%%%%%%%%%%%%%%%%%%%%%%%%%%%%%%%%%%%%%%%%
\question \textbf{An NMEA sentence.} A decoder prints
\begin{center}\small
\texttt{!AIVDM,2,1,7,B,55NQj<02D1o9L@gGG788iDF0PE8tp0000000000t:0I8840Ht>0CSPD3k2Dh,0*14}\\
\texttt{!AIVDM,2,2,7,B,00000000000,2*20}
\end{center}
\begin{parts}

	\begin{minipage}[t]{0.58\linewidth}
		\part Which channel carried it, and into how many sentences was it split?
		\begin{solution}[0.8in]
		Channel B, in 2 sentences; the 7 is the sequence number that ties them together.
		\end{solution}
	\end{minipage}\hfill%
	\begin{minipage}[t]{0.37\linewidth}
		\raggedleft \vspace{12pt}
		\ansbox[1.2in]{B, 2}
	\end{minipage}

	\begin{minipage}[t]{0.58\linewidth}
		\part Each character carries 6 bits: subtract 48 from its ASCII code (and 8 more if
		the result is above 39). The first character is \texttt{5}, ASCII 53. Find the
		message ID and say what kind of message it is.
		\begin{solution}[0.9in]
		$53 - 48 = 5$: message 5, the static and voyage report with the ship's name,
		callsign, type, size, and destination.
		\end{solution}
	\end{minipage}\hfill%
	\begin{minipage}[t]{0.37\linewidth}
		\raggedleft \vspace{12pt}
		\ansbox[1.2in]{5, name and voyage}
	\end{minipage}

	\part Explain why this message needed two sentences when a position report needs one.
		\begin{solution}[1.0in]
		A message 5 is 424 bits, which takes 71 characters at 6 bits each, more than one
		sentence holds; a position report is 168 bits, 28 characters, and fits in one. The
		same message filled two slots on the air.
		\end{solution}

\end{parts}

\newpage
%%%%%%%%%%%%%%%%%%%%%%%%%%%%%%%%%%%%%%%%%
\question \textbf{Noise and the squelch.} You run \texttt{L19-ais-rx.grc} on the lesson's
recording.
\begin{parts}

	\part You raise the Added noise slider until channel B's trace in the time plot
	swings wildly between bursts and frames start failing the CRC. Explain what has
	happened, and which slider you move to fix it.
		\begin{solution}[1.3in]
		The noise between bursts is now above the squelch threshold, so the Power Squelch
		stays open and passes noise to the FM demodulator. Symbol Sync then chases that
		noise between bursts and starts each burst out of step, so bits come out wrong and
		the CRC fails. Raising the Squelch slider above the noise closes it between bursts
		again.
		\end{solution}

	\part Explain what goes wrong if you raise the Squelch slider too far.
		\begin{solution}[0.9in]
		The squelch then shuts out the weaker ships' bursts as well as the noise, so their
		frames never reach the decoder and they drop out of the table.
		\end{solution}

	\part AIS uses a Gaussian filter with BT = 0.4. In words, say what a smaller BT
	would gain and what it would cost.
		\begin{solution}[1.1in]
		A smaller BT rounds the steps more, so the spectrum gets narrower and less power
		spills into the next channel. It also smears each bit into its neighbors, so a
		single bit's frequency swing falls short of $\pm 2.4$ kHz and noise pushes it
		across the slicer's threshold more easily.
		\end{solution}

\end{parts}

\end{questions}
